\documentclass[11p]{article}
% Packages
\usepackage{amsmath}
\usepackage{graphicx}
\usepackage{fancyheadings}
\usepackage[swedish, english]{babel}
\usepackage[
    backend=biber,
    style=authoryear-ibid,
    sorting=ynt
]{biblatex}
\usepackage[utf8]{inputenc}
\usepackage[T1]{fontenc}
\usepackage{titlesec}
\usepackage{hyperref}
%Källor
\addbibresource{references.bib}
\graphicspath{ {./images/} }

% Lite variabler
\def\email{alvin.hogdal@elev.ga.ntig.se}
\def\foottitle{Gymnasiearbete} % Kanske borde ändras
\def\name{Alvin Högdal}

\title{Gymnasiearbete \\ \small Gymnasiearbete}
\author{\name}
\date{\today}


\begin{document}



% fixar sidfot
    \lfoot{\footnotesize{\name \\ \email}}
    \rfoot{\footnotesize{\today}}
    \lhead{\sc\footnotesize\foottitle}
    \rhead{\nouppercase{\sc\footnotesize\leftmark}}
    \pagestyle{fancy}
    \renewcommand{\headrulewidth}{0.2pt}
    \renewcommand{\footrulewidth}{0.2pt}

% i Sverige har vi normalt inget indrag vid nytt stycke
    \setlength{\parindent}{0pt}
% men däremot lite mellanrum
    \setlength{\parskip}{10pt}
    \begin{otherlanguage}{swedish}


        \begin{titlepage}
            \centering

            % Title and subtitle are enclosed between two rules.
            \rule{\textwidth}{1pt}

            % Title
            \vspace{.7\baselineskip}
            {\huge \textbf{Title}}

            % Subtitle
            \vspace*{.5cm}
            {\LARGE Subtitle}

            \rule{\textwidth}{1pt}

            \vspace{1cm}

            % Set this size for the remaining titlepage.
            \large

            % Authors side by side, using two minipages as a trick.
            \begin{minipage}{.5\textwidth}
                \centering
                \name \\
                {\normalsize \url{alvinhogdal@elev.ga.ntig.se}}
            \end{minipage}%

            \vspace{3cm}

            % Report logo.
            \includegraphics[width=0.4\textwidth]{../images/NTI Gymnasiet_Symbol_print_svart.png}

            \vfill

            % University and date information at the bottom of the titlepage.
            NTI Gymnasiet Umeå \\
            Teknikprogrammet\\
            Gymnasiearbete\\
            Datum: \today \\
            Handledare: Jens Andreasson
        \end{titlepage}
        %\maketitle
        %\begin{center}
        %\includegraphics[width=0.4\textwidth]{../images/NTI Gymnasiet_Symbol_print_svart.png}
        %\end{center}
        
        
    \end{otherlanguage}


    \begin{otherlanguage}{swedish}
    \tableofcontents

    \newpage

% i Sverige har vi normalt inget indrag vid nytt stycke\setlength{\parindent}{0pt}
% men däremot lite mellanrum
\setlength{\parskip}{10pt}

\section{Inledning}

\subsection{Syfte & Frågeställning}
    Undersöka hur priset på liknande datorspel har ändrats över tid.


    Jämföra hur prisutvecklingen har varit på två liknande spel
\section{Bakgrund}

    \subsection{datorspel}
    Datorspel är
    Enligt \textcite{ComputerSweden} så defineras datorspel som något spel som kräver en dator.

    \subsection{utvecklare}
    Ett datorspel kräver att någon skapar det.
    Dessa personer eller företag som skapar ett spel har samlingsnamnet utvecklare.
    Alla som har varit med och skapat spelet ingår i namnet utvecklare.
    Hur många utvecklare ett datorspel har kan variera kraftigt beroende på hur stort spelet är.
    Det finns ingen gräns på hur många personer som kan utveckla ett spel men det krävs minst en.

    \subsection{utgivare}
    Utgivare är personerna eller företagen som ansvarar för att marknadsföra spelet.

    \subsection{Shooter}
    En \("\)shooter\("\) är en genre på datorspel.
    Att ett spel är en \("\)shooter\("\) innebär att en av huvudmekanikerna i spelet är att skjuta någon projektil på någon figur.
    Detta kan vara allt från att skjuta pilar på måltavlor till att skjuta kulor på personer.

    \subsubsection{First person shooter}
    En \("\)First person shooter\("\) eller FPS är en \("\)shooter\("\) som utgår från den spelbara karaktärens perspektiv.
    Ett exempel på en FPS är datorspelet \("\)Counter-Strike 2\("\) som är utvecklat och marknadsfört av \textcite{CounterStrike}

    \subsection{Downloadable content}
    Downloadable content eller DLC är innehåll till ett datorspel som utvecklaren och utgivaren  av någon anledning vil hålla separat.
    Det kan vara en uppdatering som ändrar någonting,

    \subsection{Battle pass}
    test
    \subsection{Micro-transactions}
    test
\section{Metod}

     Undersökningen går ut på att hitta information på när de valda spelen släpptes.
     Hur mycket det kostade då. %kanske utväcklingstid
     Vem som utvecklade spelet och vem som är utgivare.
     Alla dlc och vad de kostade när de släpptes.

    \printbibliography[title={Referenser}]

    \end{otherlanguage}

\end{document}
